\chapter{Combinatorics}\label{combinatorics}
\chapterstatus{planned}

Counting finite sets: the basic principles, the binomial and multinomial coefficients and their
identities, inclusion--exclusion, the pigeonhole principle, double counting, and the numbers that count
partitions and permutations.

\noindent\textbf{Prerequisites.} Chapter~\ref{sets} (Set Theory), Chapter~\ref{number-systems} (Number Systems) and Chapter~\ref{groups} (Groups).

\noindent\textbf{On the way to the Hodge conjecture.} This chapter is not needed on the direct
route; it belongs to the reading path of discrete mathematics
(Subsection~\ref{introduction:subsection-reading-paths}).

\noindent\textbf{Planned contents.}
\begin{enumerate}
\item The rules of sum and product, bijective proofs, counting functions, injections and subsets.
\item Binomial coefficients: Pascal's triangle, the binomial theorem, Vandermonde's identity, Lucas' theorem modulo a prime; multinomial coefficients.
\item Inclusion--exclusion: derangements, surjections, Euler's function; M\"obius inversion on partially ordered sets.
\item The pigeonhole principle and its strengthenings; Erd\H{o}s--Szekeres; Dirichlet's approximation theorem.
\item Double counting; Sperner's theorem and the LYM inequality; the Erd\H{o}s--Ko--Rado theorem.
\item Stirling numbers of both kinds, Bell numbers, Lah numbers; set partitions and cycles of permutations.
\item Integer partitions: Ferrers diagrams, conjugation, Euler's theorem on odd and distinct parts.
\item Counting under symmetry: Burnside's lemma and P\'olya enumeration.
\item Exercises with solutions.
\end{enumerate}
